\section{Desafíos y Oportunidades de Mejora}

Durante el proceso de implementación y despliegue del sistema Odoo en PERÚ FLORES TRAVEL AGENCY TOURS OPERADOR E.I.R.L., se identificaron desafíos críticos que requieren una gestión continua, así como oportunidades estratégicas para potenciar el crecimiento de la agencia en el ecosistema digital.

\subsection{Desafíos Identificados}

\begin{itemize}
    \item \textbf{Estandarización de Tarifarios Complejos:} Uno de los mayores retos fue trasladar la lógica de precios variables de la Sra. Silveria (basada en el número de pasajeros y tipos de hotel) al sistema. La configuración de estas reglas en Odoo requiere una validación precisa para asegurar que los cálculos automáticos de ganancia (30\% o 40\%) reflejen fielmente la realidad del mercado.
    \item \textbf{Transición de la Flexibilidad Manual a la Disciplina Digital:} La agencia operaba con una alta dependencia de la memoria y la "flexibilidad manual" de herramientas como WhatsApp y cuadernos. El desafío principal radica en adoptar la disciplina de registro sistemático de cada oportunidad comercial y reserva en el CRM para evitar que la información vuelva a quedar aislada en chats personales.
    \item \textbf{Dependencia de la Conectividad en Zonas Turísticas:} Dado que Odoo es una plataforma 100\% web, la inestabilidad del internet en ciertas zonas de operación o durante traslados en carretera representa una limitante operativa. Se requiere capacitar al personal en el uso eficiente de la aplicación móvil para consultas rápidas de itinerarios en entornos con señal limitada.
    \item \textbf{Gestión del Cambio en el Capital Humano:} El personal posee una alta experiencia turística pero una formación digital empírica. Superar la resistencia al cambio y asegurar que tanto la gerencia como el equipo operativo utilicen el ERP como su herramienta principal de trabajo es crítico para el éxito a largo plazo.
\end{itemize}

\subsection{Oportunidades de Mejora}

\begin{itemize}
    \item \textbf{Fortalecimiento de la Reputación mediante Feedback Automatizado:} Existe una gran oportunidad para capturar el "boca a boca" digital mediante el módulo de Encuestas. Al automatizar el envío de cuestionarios post-viaje, la agencia puede alimentar su sitio web con testimonios reales, mejorando la confianza de nuevos clientes.
    \item \textbf{Implementación de Pagos en Línea:} Actualmente, la agencia depende de depósitos por Western Union o bancos, lo cual es fastidioso para los clientes internacionales. Integrar pasarelas de pago (PayPal o similares) en el flujo de ventas de Odoo permitiría cerrar ventas de manera inmediata y segura.
    \item \textbf{Uso de Inteligencia Artificial para la Comunicación Multilingüe:} La agencia atiende principalmente mercados brasileños y de habla inglesa. El uso de las herramientas de IA integradas en Odoo para mejorar la redacción y traducción automática de cotizaciones permite profesionalizar la comunicación sin depender de traductores externos.
    \item \textbf{Analítica de Datos para la Toma de Decisiones:} El sistema genera información valiosa sobre rentabilidad por paquete y tasas de conversión. Consolidar el uso de Dashboards de BI permitirá a la gerencia identificar qué destinos (Cusco, Arequipa o Puno) son más rentables y ajustar su estrategia de marketing en consecuencia.
\end{itemize}
